\scanpage{191}
définie par
\[
 \tau\longrightarrow\tau'=\Sigma-\tau,
\]

on aura
\[
 \tau'=\operatorname{Tr}yxz=\operatorname{Tr}xzy=\operatorname{Tr}zyx.
\]

Les \unsure{élémts|éléments} du groupe symétrique $\unsure{\mathfrak S_3|\mathscr S_3}$ échangent entre eux les neuf générateurs $\underline b,\underline c,\underline t$, et aussi $\tau$ et $\tau'$ (les permutations circ. laissent $\tau,\tau'$ fixes, les transpositions échangent $\tau,\unsure{\tau''|\tau'}$).

En \deleted{termes} tant qu'anneau où $\mathfrak S_3$ opère, on peut décrire $B$ comme
% The curved transposition mark leads from the group of t's to the point after the double-primed c.
\[
 B=k[\underline b,\underline b',\underline b'',\underline c,\underline c',\underline c'',\inserted{\underline t,\underline t',\underline t'',}\underline\tau,\underline\tau']/
 \bigl(\underset{\substack{\text{équations séparé-}\\\text{ment invariantes}\\\text{sous }\mathfrak S_3}}{\underbrace{\underline\tau+\underline\tau'-\Sigma},\ \underbrace{\tau\tau'-\Pi}}\bigr)
\]

avec l'opération évidente de $\mathfrak S_3$ sur l'anneau des polynômes à onze variables $\underline b,\underline{\unsure{C|c}},\underline t,\underline\tau,\underline\tau'$.

\uline{NB} Quand $2$ est inv. dans $k$, et dans ce
% archive sheet number: 190
